\renewcommand{\tema}{Examen Diagnóstico - Ondas}

% -------------------------------------------------------
%  PREGUNTAS
% -------------------------------------------------------

\twocolumn
\section*{PREGUNTAS}

\begin{preguntabox}
\subsection*{Pregunta 1}
¿Cuál de los siguientes es un ejemplo de onda \textbf{transversal}?

\begin{enumerate}
    \item[A)] El sonido que viaja por el aire.
    \item[B)] La onda que recorre una cuerda de guitarra al pulsarla.
    \item[C)] Una onda sísmica de tipo P que viaja por la corteza terrestre.
    \item[D)] Las ondas de presión generadas por una explosión.
\end{enumerate}
\end{preguntabox}

\begin{preguntabox}
\subsection*{Pregunta 2}
En una onda longitudinal, las partículas del medio vibran:

\begin{enumerate}
    \item[A)] Perpendicularmente a la dirección de propagación.
    \item[B)] En círculos alrededor del eje de propagación.
    \item[C)] Paralelamente a la dirección de propagación.
    \item[D)] En dirección aleatoria, independiente de la propagación.
\end{enumerate}
\end{preguntabox}

\begin{preguntabox}
\subsection*{Pregunta 3}
Una onda tiene una frecuencia de 5 Hz. ¿Cuál es su periodo?

\begin{enumerate}
    \item[A)] 5 s
    \item[B)] 0.5 s
    \item[C)] 0.2 s
    \item[D)] 2 s
\end{enumerate}
\end{preguntabox}

\begin{preguntabox}
\subsection*{Pregunta 4}
Una onda tiene una longitud de onda de 4 m y una frecuencia de 25 Hz. ¿Cuál es su velocidad de propagación?

\begin{enumerate}
    \item[A)] 6.25 m/s
    \item[B)] 29 m/s
    \item[C)] 21 m/s
    \item[D)] 100 m/s
\end{enumerate}
\end{preguntabox}

\begin{preguntabox}
\subsection*{Pregunta 5}
Una onda pasa de un medio donde viaja a 300 m/s a otro donde viaja a 150 m/s. Si la frecuencia de la onda es de 50 Hz, ¿cuál es la longitud de onda en el segundo medio?

\begin{enumerate}
    \item[A)] 6 m
    \item[B)] 3 m
    \item[C)] 1 m
    \item[D)] 12 m
\end{enumerate}
\end{preguntabox}

\begin{preguntabox}
\subsection*{Pregunta 6}
¿En cuál de las siguientes situaciones la difracción es más significativa?

\begin{enumerate}
    \item[A)] Una onda de radio AM ($\lambda \approx 300$ m) que encuentra un edificio de 20 m de ancho.
    \item[B)] Un rayo de luz visible ($\lambda \approx 500$ nm) que pasa por una ventana de 1 m de ancho.
    \item[C)] El sonido audible ($\lambda \approx 1$ m) que pasa por una puerta de 1 m de ancho.
    \item[D)] Una microonda ($\lambda \approx 3$ cm) que incide sobre una pared de 10 m de ancho.
\end{enumerate}
\end{preguntabox}

\begin{preguntabox}
\subsection*{Pregunta 7}
Dos ondas de igual amplitud $A = 3$ cm se superponen en fase (cresta con cresta). ¿Cuál es la amplitud de la onda resultante?

\begin{enumerate}
    \item[A)] 0 cm
    \item[B)] 3 cm
    \item[C)] 6 cm
    \item[D)] 1.5 cm
\end{enumerate}
\end{preguntabox}

\begin{preguntabox}
\subsection*{Pregunta 8}
Una onda se propaga en un medio con velocidad $v$ y longitud de onda $\lambda$. Si la longitud de onda se duplica y la velocidad permanece constante, la nueva frecuencia es:

\begin{enumerate}
    \item[A)] $4f$
    \item[B)] $2f$
    \item[C)] $f$
    \item[D)] $\dfrac{f}{2}$
\end{enumerate}
\end{preguntabox}

\begin{preguntabox}
\subsection*{Pregunta 9}
Una onda viaja a 60 m/s y tiene un periodo de 0.3 s. ¿Cuál es su longitud de onda?

\begin{enumerate}
    \item[A)] 200 m
    \item[B)] 0.005 m
    \item[C)] 18 m
    \item[D)] 60.3 m
\end{enumerate}
\end{preguntabox}

\begin{preguntabox}
\subsection*{Pregunta 10}
Un barco emite una señal de sonar con longitud de onda de 50 m y frecuencia de 4 Hz. La señal rebota en el fondo del mar y regresa al barco 6 segundos después. ¿Cuál es la profundidad del mar en ese punto?

\begin{enumerate}
    \item[A)] 1 200 m
    \item[B)] 600 m
    \item[C)] 300 m
    \item[D)] 2 400 m
\end{enumerate}
\end{preguntabox}

% -------------------------------------------------------
%  RESPUESTAS Y RETROALIMENTACIÓN
% -------------------------------------------------------

\newpage
\onecolumn
\section*{RESPUESTAS Y RETROALIMENTACIÓN}

\begin{respuestabox}
\subsection*{Pregunta 1}
\textcolor{verdecorrecto}{\textbf{Respuesta correcta: B) La onda que recorre una cuerda de guitarra al pulsarla.}}

\textbf{Justificación:}\\
En una onda transversal, las partículas vibran \textit{perpendicularmente} a la dirección de propagación. Al pulsarla, la cuerda sube y baja (vibración vertical) mientras la onda avanza horizontalmente a lo largo de la cuerda.

\textbf{Respuestas incorrectas:}
\begin{itemize}
    \item \textbf{A) Sonido en el aire:} El sonido es longitudinal — las moléculas de aire se comprimen y rarifican en la misma dirección en que viaja el sonido.
    \item \textbf{C) Onda sísmica tipo P:} Las ondas P (primarias) son longitudinales. Las ondas S (secundarias) sí son transversales.
    \item \textbf{D) Ondas de presión:} Las variaciones de presión implican compresión y rarefacción en la misma dirección de propagación: son longitudinales.
\end{itemize}
\end{respuestabox}

\begin{respuestabox}
\subsection*{Pregunta 2}
\textcolor{verdecorrecto}{\textbf{Respuesta correcta: C) Paralelamente a la dirección de propagación.}}

\textbf{Justificación:}\\
Por definición, en una onda longitudinal la perturbación (compresión/rarefacción) ocurre en la misma dirección en que se propaga la onda. El sonido en el aire es el ejemplo más común.

\textbf{Respuestas incorrectas:}
\begin{itemize}
    \item \textbf{A) Perpendicularmente:} Esa es la definición de onda \textit{transversal}, no longitudinal.
    \item \textbf{B) En círculos:} Las partículas de agua en la superficie describen trayectorias elípticas, lo cual corresponde a ondas superficiales, no a ondas longitudinales puras.
    \item \textbf{D) Dirección aleatoria:} Ningún tipo de onda tiene vibración aleatoria; la dirección de oscilación está determinada por el tipo de onda.
\end{itemize}
\end{respuestabox}

\begin{respuestabox}
\subsection*{Pregunta 3}
\textcolor{verdecorrecto}{\textbf{Respuesta correcta: C) 0.2 s}}

\textbf{Justificación:}\\
\textbf{Fórmula:} $T = \dfrac{1}{f}$

$$T = \frac{1}{5 \text{ Hz}} = 0.2 \text{ s}$$

\textbf{Respuestas incorrectas:}
\begin{itemize}
    \item \textbf{A) 5 s:} Error de confundir $T$ con $f$; el alumno simplemente copió el valor de la frecuencia.
    \item \textbf{B) 0.5 s:} Error de dividir $\frac{1}{2}$ en lugar de $\frac{1}{5}$; confusión con $f = 2$ Hz.
    \item \textbf{D) 2 s:} Error de calcular $\frac{f}{?}$ incorrectamente, o confundir $T = f$ con $T = \frac{f}{10}$.
\end{itemize}
\end{respuestabox}

\begin{respuestabox}
\subsection*{Pregunta 4}
\textcolor{verdecorrecto}{\textbf{Respuesta correcta: D) 100 m/s}}

\textbf{Justificación:}\\
\textbf{Fórmula:} $v = \lambda \cdot f$

$$v = 4 \text{ m} \times 25 \text{ Hz} = 100 \text{ m/s}$$

\textbf{Respuestas incorrectas:}
\begin{itemize}
    \item \textbf{A) 6.25 m/s:} Error de dividir $\lambda \div f = 4 \div 25$ en lugar de multiplicar.
    \item \textbf{B) 29 m/s:} Error de sumar $\lambda + f = 4 + 25$; confunde la operación.
    \item \textbf{C) 21 m/s:} Error de restar $25 - 4$; otra confusión de operación.
\end{itemize}
\end{respuestabox}

\begin{respuestabox}
\subsection*{Pregunta 5}
\textcolor{verdecorrecto}{\textbf{Respuesta correcta: B) 3 m}}

\textbf{Justificación:}\\
Al refractarse, la frecuencia \textbf{no cambia}. En el segundo medio se usa la nueva velocidad:
$$\lambda_2 = \frac{v_2}{f} = \frac{150 \text{ m/s}}{50 \text{ Hz}} = 3 \text{ m}$$

\textbf{Respuestas incorrectas:}
\begin{itemize}
    \item \textbf{A) 6 m:} Error de usar la velocidad del primer medio: $300/50 = 6$ m. Olvida que al cambiar de medio, la velocidad cambia.
    \item \textbf{C) 1 m:} Error de calcular $v_2 \div v_1 = 150 \div 300 = 0.5$, luego dividirlo entre $f$ incorrectamente.
    \item \textbf{D) 12 m:} Error de sumar las dos velocidades: $(300 + 150)/50 = 9$, o de multiplicar $v_2 \times f$ en lugar de dividir.
\end{itemize}
\end{respuestabox}

\begin{respuestabox}
\subsection*{Pregunta 6}
\textcolor{verdecorrecto}{\textbf{Respuesta correcta: C) El sonido audible que pasa por una puerta de 1 m de ancho.}}

\textbf{Justificación:}\\
La difracción es significativa cuando el tamaño de la apertura es \textbf{comparable} a la longitud de onda ($d \approx \lambda$):
$$d = 1 \text{ m}, \quad \lambda \approx 1 \text{ m} \quad \Rightarrow \quad d \approx \lambda \quad \checkmark$$

Por eso escuchas lo que ocurre al otro lado de una puerta aunque no tengas línea de vista directa.

\textbf{Respuestas incorrectas:}
\begin{itemize}
    \item \textbf{A) Radio AM y edificio:} $\lambda = 300$ m $\gg$ $d = 20$ m. La onda sí se difracta, pero la condición $d \approx \lambda$ no se cumple: el obstáculo es mucho más pequeño que la longitud de onda.
    \item \textbf{B) Luz visible y ventana:} $\lambda \approx 5 \times 10^{-7}$ m $\ll$ $d = 1$ m. La apertura es millones de veces más grande que $\lambda$: la luz viaja en línea recta sin difractarse apreciablemente.
    \item \textbf{D) Microondas y pared:} $\lambda = 3$ cm $\ll$ $d = 10$ m. La pared es cientos de veces más grande que $\lambda$: difracción despreciable.
\end{itemize}
\end{respuestabox}

\begin{respuestabox}
\subsection*{Pregunta 7}
\textcolor{verdecorrecto}{\textbf{Respuesta correcta: C) 6 cm}}

\textbf{Justificación:}\\
\textbf{Interferencia constructiva:} $A_{\text{total}} = A_1 + A_2$

$$A_{\text{total}} = 3 \text{ cm} + 3 \text{ cm} = 6 \text{ cm}$$

Cuando dos ondas llegan en fase (cresta con cresta), sus amplitudes se suman.

\textbf{Respuestas incorrectas:}
\begin{itemize}
    \item \textbf{A) 0 cm:} Corresponde a interferencia \textit{destructiva} con $A_1 = A_2$, no constructiva.
    \item \textbf{B) 3 cm:} Error de pensar que la amplitud no cambia al superponerse; ignora el principio de superposición.
    \item \textbf{D) 1.5 cm:} Error de dividir la amplitud entre 2; confunde promedio con superposición.
\end{itemize}
\end{respuestabox}

\begin{respuestabox}
\subsection*{Pregunta 8}
\textcolor{verdecorrecto}{\textbf{Respuesta correcta: D) $\dfrac{f}{2}$}}

\textbf{Justificación:}\\
De $v = \lambda \cdot f$, con $v$ constante: $f = \dfrac{v}{\lambda}$.

Si $\lambda' = 2\lambda$:
$$f' = \frac{v}{2\lambda} = \frac{f}{2}$$

Al duplicar la longitud de onda con velocidad constante, la frecuencia se reduce a la mitad. La relación entre $\lambda$ y $f$ es inversamente proporcional.

\textbf{Respuestas incorrectas:}
\begin{itemize}
    \item \textbf{A) $4f$:} Confunde la relación inversa con una relación cuadrática.
    \item \textbf{B) $2f$:} Error de pensar que si $\lambda$ aumenta, $f$ también aumenta (confunde proporcionalidad directa con inversa).
    \item \textbf{C) $f$:} Error de pensar que la frecuencia no cambia aunque cambie la longitud de onda con $v$ fija.
\end{itemize}
\end{respuestabox}

\begin{respuestabox}
\subsection*{Pregunta 9}
\textcolor{verdecorrecto}{\textbf{Respuesta correcta: C) 18 m}}

\textbf{Justificación:}\\
\textbf{Paso 1:} Usar $v = \lambda / T$, despejando $\lambda$:
$$\lambda = v \cdot T = 60 \text{ m/s} \times 0.3 \text{ s} = 18 \text{ m}$$

\textbf{Respuestas incorrectas:}
\begin{itemize}
    \item \textbf{A) 200 m:} Error de calcular $v \div T = 60 \div 0.3 = 200$; invierte la fórmula $\lambda = v \cdot T$.
    \item \textbf{B) 0.005 m:} Error de calcular $T \div v = 0.3 \div 60$; invierte todos los términos.
    \item \textbf{D) 60.3 m:} Error de sumar $v + T = 60 + 0.3$; confunde la operación.
\end{itemize}
\end{respuestabox}

\begin{respuestabox}
\subsection*{Pregunta 10}
\textcolor{verdecorrecto}{\textbf{Respuesta correcta: B) 600 m}}

\textbf{Justificación:}\\
\textbf{Paso 1:} Calcular la velocidad del sonido en el agua:
$$v = \lambda \cdot f = 50 \text{ m} \times 4 \text{ Hz} = 200 \text{ m/s}$$

\textbf{Paso 2:} La señal recorre el camino de ida y vuelta en 6 s, por lo que el tiempo de ida es $t = 3$ s:
$$d = v \cdot t = 200 \text{ m/s} \times 3 \text{ s} = 600 \text{ m}$$

\textbf{Respuestas incorrectas:}
\begin{itemize}
    \item \textbf{A) 1 200 m:} Error de no dividir el tiempo entre 2; olvida que la señal hace el recorrido de ida y vuelta.
    \item \textbf{C) 300 m:} Error de usar $t = 6$ s pero con $v = 100$ m/s (calcula $v = \lambda + f$ en lugar de $v = \lambda \cdot f$).
    \item \textbf{D) 2 400 m:} Error de multiplicar $v \times t$ sin dividir el tiempo: $200 \times 6 \times 2$.
\end{itemize}
\end{respuestabox}